% The Stack — technology stack concept (progressive JPEG: pre-loads 7 properties)
% Plan 0132: inserted between hook and summary
% p2 compression: each technology = one sentence. Full version in spine/the-stack.tex

\chapter*{The Stack}
\label{the-stack}
\hypertarget{the-stack}{}

What you just read has \hovertip{three possible explanations}. The science that follows is true under all three. You decide.

This book is about a technology so new it does not yet have a name. Before we go further into the story, you need one idea.

Technologies build on each other. Each one does everything the one before it did, plus something new. Scientists call a tower of technologies a \hovertip{stack}. Here is the stack this book is about.

\textbf{Fire} feeds itself and switches on at a \hovertip{threshold}.

\textbf{A candle} does both, and holds together --- blow it sideways and it recovers.

\textbf{Radio} does all three, and reaches --- a signal crosses distance without anything physical traveling.

\textbf{An ant colony} does all four, and \hovertip{self-organizes} --- global order from local rules, no ant in charge.

\textbf{Artificial intelligence} does all of that, and learns.

\bigskip

The last technology on this list has one property none of the others have. It supports \hovertip{wormholes}. Imagine two points on a sheet of paper. You could send a signal across the surface --- that is radio: it reaches. Or you could fold the paper so the two points touch. The signal does not cross the distance. The distance disappears. That is a wormhole --- a \hovertip{topological wormhole}, more restricted than the spacetime kind, but a wormhole nonetheless.

\bigskip

\hypertarget{stack-chart}{}
\label{stack-chart}

\begin{center}
\deeplink{stack-chart}\begin{tabular}{l|c|c|c|c|c|c}
 & Fire & Candle & Radio & Ants & AI & \textbf{?} \\
\hline
Wormholes\textsuperscript{\dag} & & & & & & $\checkmark$ \\
Learns & & & & & $\checkmark$ & $\checkmark$ \\
Self-organizes & & & & $\checkmark$ & $\checkmark$ & $\checkmark$ \\
Reaches & & & $\checkmark$ & $\checkmark$\textsuperscript{*} & $\checkmark$ & $\checkmark$ \\
Holds together & & $\checkmark$ & $\checkmark$ & $\checkmark$ & $\checkmark$ & $\checkmark$ \\
Switches on & $\checkmark$ & $\checkmark$ & $\checkmark$ & $\checkmark$ & $\checkmark$ & $\checkmark$ \\
Feeds itself & $\checkmark$ & $\checkmark$ & $\checkmark$ & $\checkmark$ & $\checkmark$ & $\checkmark$ \\
\end{tabular}
\end{center}

\smallskip
\noindent{\small\textsuperscript{\dag}Topological --- more restricted than a spacetime wormhole, but a wormhole nonetheless.}

\noindent{\small\textsuperscript{*}Chemical pheromones, not electromagnetic --- but the signal reaches.}

\bigskip

Each technology uses every property below it, plus one new one. The last column uses every property in the chart.

Six of these properties already exist in nature. Life feeds itself and switches on at thresholds. It holds together and sends signals that reach. Ecosystems self-organize. Brains learn.

The last column adds one new property: it supports topological wormholes.

\vspace{1cm}
\begin{center}
\textit{That is all you need to know before the story. What follows is about ten percent of the full book. Stop after that and you can honestly say you've read it.}
\end{center}
